\documentclass[11pt,a4paper]{article}
\usepackage{../shared/hanzocover}
\usepackage[utf8]{inputenc}
\usepackage[T1]{fontenc}
% Shared preamble for the compute-market series.
% One style, one vocabulary, inputted by every paper so the series reads as one voice.
\usepackage[margin=1in]{geometry}
\usepackage{booktabs}
\usepackage{amsmath}
\usepackage{amssymb}
\usepackage{amsthm}
\usepackage{graphicx}
\usepackage{xcolor}
\usepackage{hyperref}
\hypersetup{colorlinks=true, linkcolor=blue, urlcolor=blue, citecolor=blue}
\usepackage{enumitem}
\usepackage{listings}
\usepackage{array}
\usepackage{tabularx}
\usepackage{siunitx}
\sisetup{group-separator={,}, group-minimum-digits=4}

\definecolor{kw}{rgb}{0.13,0.29,0.53}
\definecolor{cmt}{rgb}{0.40,0.40,0.40}
\definecolor{str}{rgb}{0.16,0.50,0.16}
\lstset{
  basicstyle=\ttfamily\footnotesize,
  keywordstyle=\color{kw}\bfseries,
  commentstyle=\color{cmt}\itshape,
  stringstyle=\color{str},
  showstringspaces=false,
  breaklines=true,
  columns=fullflexible,
  frame=single,
  framesep=4pt,
}

\newtheorem{definition}{Definition}
\newtheorem{proposition}{Proposition}
\newtheorem{remark}{Remark}

\newcommand{\code}[1]{\texttt{#1}}
\newcommand{\repo}[1]{\texttt{#1}}

% Series vocabulary — used identically in every paper.
\newcommand{\job}{\ensuremath{J}}
\newcommand{\bundle}{\ensuremath{R}}
\newcommand{\cost}{\ensuremath{\mathcal{C}}}

\tolerance=800
\emergencystretch=3em


\title{Las Vegas Workloads:\\Verifiable Search in a Compute Market}

\author{
Hanzo AI Research\thanks{Corresponding author: research@hanzo.ai} \\
\textit{Hanzo Industries Inc (Techstars '17)} \\
Los Angeles, California \\
\texttt{https://hanzo.ai}
}

\date{2026}

\begin{document}
\hanzocoverpage

\begin{abstract}
The obvious first workload for a decentralized compute market is inference,
because inference is what people want to buy. It is also close to the worst
workload the design could start with, because there is no cheap way to check it.

Verifiable search is the opposite. A search job is enormous to compute and
trivial to check, the check is binary, an invalid answer simply fails rather
than requiring adjudication, the work is embarrassingly parallel across
heterogeneous supply, and the input is public so provider selection carries no
privacy constraint. Fraud is not detected --- it is impossible, which removes an
entire layer of protocol.

We formalise the class, contrast its market properties with inference, and work
through a live example: hash-based signature search for quantum-resistant
Bitcoin spends. The example is instructive beyond itself, because analysing it
produces a design rule that generalises to every proof-of-work-flavoured
construction, and because the block-space arithmetic shows how a scheme's
container choice can dominate its cryptography.
\end{abstract}

\tableofcontents
\newpage

\section{The class}

\begin{definition}[Las Vegas workload]
A job is a pair $(\mathcal{V}, p)$ where $\mathcal{V}(x, w) \in \{0,1\}$ is a
predicate computable in time negligible against the search, and each independent
trial succeeds with probability $p$. The provider's task is to find any $w$ with
$\mathcal{V}(x,w) = 1$. Expected work is $1/p$ trials; verification cost is
independent of the work performed.
\end{definition}

Four market properties follow directly, and together they are unusual enough to
be worth listing.

\paragraph{Verification is free.} $\mathcal{V}$ is by construction cheap. No
proof system, no replication, no committee. The job sits at tier T1 of the
verification ladder at zero marginal cost.

\paragraph{Fraud is impossible rather than detectable.} A provider cannot submit
a wrong answer that passes, because passing is the definition of right. There is
nothing to slash for, because there is nothing a dishonest provider can gain.
The worst available misbehaviour is doing nothing, which is a liveness problem
and is priced by the deadline.

\paragraph{Sharding is free.} The search space partitions cleanly. Ten thousand
providers can work disjoint nonce ranges with no coordination beyond the initial
assignment, and heterogeneous hardware simply contributes at whatever rate it
manages.

\paragraph{The input is public.} There is no confidential material to protect,
so provider selection is unconstrained by privacy and the job can be broadcast
openly. Private preimages, where the construction has them, never leave the
owner's device --- only the public search is outsourced.

\section{Why inference is harder}

Inference has no cheap $\mathcal{V}$. Given a token stream, there is no
sub-linear test that it came from the model that was paid for. Options are
recomputation (as expensive as the job), replication (multiplies the bill),
algebraic spot-checking of the matrix products (good, but silent about the
nonlinearities and about which weights were used), or a succinct proof (three to
six orders of magnitude of prover overhead).

This is not an argument against building an inference market. It is an argument
for not making the verification layer's first load-bearing test be its hardest
case. A market that carries search volume from day one has real verified
throughput while its inference verification story matures.

\section{A worked example}

\subsection{The construction}

Consider a scheme that makes a Bitcoin spend's security rest on hash preimage
resistance rather than on the discrete logarithm, using only opcodes already
deployed --- so that coins can be protected against a Shor-capable adversary
without waiting for a consensus change.

The mechanism is forced by a limit. Bitcoin's legacy script cannot compute the
index-selection function of a hash-based one-time signature within its
201-opcode budget. So instead of computing which key indices a message selects,
the script hardcodes the index set and the signer grinds the transaction ---
varying a nonce field --- until its hash lands on that set. The grind is the
mechanism, and it is a Las Vegas workload in exactly the sense defined above.

\subsection{The arithmetic}

For a scheme committing to 9 indices from 150 in a first round and 8 from the
remaining 141 in a second:

\begin{table}[h]
\centering
\begin{tabular}{lrr}
\toprule
Quantity & Value & Bits \\
\midrule
$\binom{150}{9}$ & \num{8.29e13} & $2^{46.2}$ \\
$\binom{141}{8}$ & \num{3.17e12} & $2^{41.5}$ \\
Sum, rounds sequential & \num{8.32e13} & $2^{46.2}$ \\
Product, rounds joint & \num{2.63e26} & $2^{87.8}$ \\
\bottomrule
\end{tabular}
\caption{Search size under the two possible round structures.}
\end{table}

Whether the rounds compose additively or multiplicatively is not a detail. At
$1.5 \times 10^{10}$ hashes per second per accelerator, the additive reading is
about \SI{1.5}{\hour} on a single device; the multiplicative reading is
$5.6 \times 10^{8}$ device-years. Those are not two versions of the same product.
One requires no marketplace at all; the other cannot be completed by any
marketplace.

\subsection{The design rule}

The more interesting observation concerns the adversary. In a grind-secured
scheme the honest signer's work and the attacker's work are the same
computation: once a spend is broadcast, its preimages are public, and an
adversary wishing to redirect the output must complete an identical search
before the transaction confirms. Security is therefore
\begin{equation}
\text{margin} \;=\; \frac{\text{honest precomputation time}}{\text{confirmation window}} \cdot \frac{\text{honest rate}}{\text{adversary rate}} ,
\end{equation}
and the second factor is where such schemes die. Against an adversary holding
SHA-256 application-specific silicon, a $10^{21}$ grind takes a mid-size mining
pool about \SI{1000}{\second} and the whole network about \SI{1}{\second}, while
ten thousand accelerators need \num{77} days.

\begin{proposition}[Grind selection rule]
Never grind a primitive for which your adversary owns dedicated silicon.
\end{proposition}

Applied here: if the ground function is a double SHA-256 the construction is
unsound in practice regardless of its combinatorics, because the entire Bitcoin
mining industry is an oracle for exactly that function. If it is
RIPEMD-160 composed with SHA-256 --- the \code{HASH160} opcode --- no deployed
miner accelerates it, general-purpose accelerators are the fastest available
hardware, and the symmetry breaks in the defender's favour. One opcode decides
whether the scheme works.

\subsection{Container arithmetic}

A second lesson is that where a script lives can dominate what it does.

\begin{table}[h]
\centering
\begin{tabular}{lrr}
\toprule
 & Bare script & Committed script path \\
\midrule
Funding output & \num{9661} vB & \num{43} vB \\
Unspent-output-set cost until spent & \SI{9661}{\byte} & \SI{43}{\byte} \\
Spend & \num{646} vB & \num{2645} vB \\
Lifecycle total & \num{10318} vB & \num{2699} vB \\
One million protected addresses & \SI{9.66}{\giga\byte} & \SI{43}{\mega\byte} \\
\bottomrule
\end{tabular}
\caption{Block-space cost of a \SI{9650}{\byte} script by container, for one
input and one output revealing seventeen 32-byte preimages.}
\end{table}

The inversion is instructive. The bare-script \emph{spend} is cheaper, because
the script was already paid for when the output was created. That is precisely
the problem: this construction is insurance, and the bare container charges the
entire premium on day one, on chain, permanently, whether or not the policy is
ever claimed. A committed script path defers about ninety-eight percent of the
cost to the claim and reduces the permanent burden on every validating node by a
factor of \num{225}.

There is a caveat that outranks the table. If the construction depends on a
script-evaluation quirk that exists only in legacy script and was removed by the
segregated-witness upgrade, then it cannot be relocated at all, and the size
ceiling, the opcode ceiling and the grind stop being implementation choices and
become structural. That is the first thing to establish about any such design,
because it determines whether the remaining questions are about tuning or about
redesign.

\subsection{The construction that needs no search}

Worth stating because it reframes the whole exercise: the grind exists only
because legacy script could not compute the index selection. The tapscript
upgrade removed both the opcode limit and the script size limit, and hash-chain
one-time signatures have since been implemented in that environment as bit
commitments. Such a signature commits to the message digest directly. There is
no index selection, therefore no grind, therefore no search job, therefore no
compute market involvement --- and the construction is studied rather than
bespoke.

A market should want to know this. Selling compute for a search that a better
construction eliminates is not a business; recognising it and pricing the
remaining genuine search workloads is.

\section{Market mechanics for search}

\paragraph{Job specification.} Declare $\mathcal{V}$ as content-addressed code,
the success probability $p$, the search domain, and a deadline. Expected cost is
$1/p$ trials divided by the aggregate rate purchased.

\paragraph{Partitioning.} Assign disjoint domain ranges. A provider can lie about
which range it searched, but cannot fabricate a solution, so the worst outcome is
duplicated effort rather than theft. Overlapping assignments with a modest
redundancy factor are usually cheaper than the coordination needed to prevent
overlap.

\paragraph{Settlement.} Winner-take-all invites providers to hoard a solution
until the last moment when a race is close, and pays nothing to the participants
who bounded the risk. Proportional payment on assigned-and-attested range
coverage plus a finder's bonus is better behaved, at the cost of needing a weak
attestation of coverage --- which is one of the few places in this class where
some trust re-enters.

\paragraph{Pricing.} Expected work is known in advance, which makes search one of
the few compute workloads that can be quoted as a fixed price with a genuine
distribution rather than a guess. Variance is geometric with parameter $p$; the
market can offer a fixed-price contract and absorb the variance across jobs,
which is an insurance product and a real revenue line.

\section{Summary}

Search is the workload where a decentralized market's verification story is
free, its adjudication story is empty, and its sharding story is trivial.
Starting there is not a compromise; it is the case the architecture handles
perfectly, and it buys the time needed to do inference verification properly.

The example examined here also carries a warning worth generalising: a search
workload created by a design limitation may evaporate when the limitation is
lifted. The durable search workloads are the ones where the search is intrinsic
to the problem, not an artifact of the container.

\end{document}
